\chapter*{Annexes}
\addcontentsline{toc}{chapter}{Annexes}

\section*{Annexe A : Description du Dataset}
Le dataset utilisé dans ce projet est le jeu de données IBM HR Analytics Employee Attrition \& Performance disponible sur Kaggle.  
Il contient des informations sur les employés telles que :
\begin{itemize}
    \item Âge
    \item Département
    \item Salaire mensuel
    \item Niveau d’éducation
    \item Satisfaction au travail
    \item Nombre d’années d’expérience
    \item Taux d’attrition
\end{itemize}

\section*{Annexe B : Technologies Utilisées}
Les principales technologies utilisées dans ce projet sont :
\begin{itemize}
    \item Python
    \item Pandas
    \item NumPy
    \item SQL Server
    \item Power BI
    \item Jupyter Notebook
\end{itemize}

\section*{Annexe C : Exemple de Requête SQL}
\begin{verbatim}
SELECT Department,
       AVG(MonthlyIncome) AS AverageSalary
FROM Employees
GROUP BY Department;
\end{verbatim}

\section*{Annexe D : Captures d’Écran}
Cette section contient :
\begin{itemize}
    \item Les tableaux de bord Power BI
    \item Les visualisations analytiques
    \item Les résultats des analyses RH
\end{itemize}

\section*{Annexe E : Objectifs du Projet}
Les principaux objectifs du projet sont :
\begin{itemize}
    \item Construire un Data Warehouse RH
    \item Analyser les performances des employés
    \item Étudier les facteurs d’attrition
    \item Faciliter la prise de décision grâce à la Business Intelligence
\end{itemize}